\chapter[Basic Musical Terms and Concepts \textsuperscript{Novotny}]{Basic Musical Terms and Concepts}
\label{chap:Basic Musical Terms}

\paragraph{Organisation of the Chapter}
This chapter explains basic terms and concepts related to music, music theory, and musical instruments that will be used throughout the following chapters. Note that the treatment of the concepts in this chapter is very brief and informal. It is inteded to give the layman an overview of the topic that is sufficient for understanding the following chapters.

\section{Musical Terms}
\paragraph{Note}
\index{Note}
Primarily a note consits of:

\begin{itemize}
	\item{pitch}
	\item{duration}
\end{itemize}

In the context of this project a note contains the following additional information:
\begin{itemize}
	\item{position within the score (see \ref{sec:score})}
	\item{where on the fretboard the note is played (relevant for string instruments)}
	\item{the way the note is played (in terms of playing techniques)}
\end{itemize}

\paragraph{Chord}
\index{Chord}
\index{Triad}
A chord is a set of at least three notes played at the same time. 
The term \textit{Triad} refers to a chord that consists of exactly three notes. Triads are the only types of chords relevant to the project. This does not mean that chords of more than three notes contained in scores can not be displayed, recognised and played properly, refer to later chapters for clarification.

There are four types of triads, commonly referred to as chord qualities:
\begin{itemize}
	\item{major}
	\item{minor}
	\item{augmented}
	\item{diminished}
\end{itemize}

A chord is typically referred to by its root note and its quality (for example \textit{D Major, E Minor}).

A chord is called \textit{inverted} if the lowest note of the chord is not the root note. A \textit{C Major} chord consists of the notes \textit{C, E, G}. If a \textit{C} chord is rearranged in a way such that \textit{E} is the lowest note of the chord this \textit{inversion} of the chord is commonly referred to as a \textit{C/E} chord.

\paragraph{Score}
\label{sec:score}
Scores are analog, visual representations of musical pieces. A score consists of notes and are notated in one of the forms of musical notations discussed in the next section.

\section{Musical Notation}
\paragraph{Piano Roll}
\index{Piano Roll}
Piano Rolls were developed as a storage medium for musical performances primarily used for operating self-playing pianos. Piano Rolls are continous rolls of papers with perforations representing single notes. The horizontal position of the perforation represents the pitch of the note, the length of the perforation represents the duration.

Throughout the last years \textit{Digital Audio Workstations} have adopted this concept for displaying MIDI scores.
The following graphic shows the notation of a melody as represented by the popular open source score notation program \textit{TuxGuitar}.

\includegraphics[width=\textwidth]{pianoroll_example}

\paragraph{Modern Western Music Notation (Standard Music Notation, Score)}
\index{Score}
\index{Standard Music Notation}
\index{Modern Western Music Notation}

Modern western music notation is based on five lines representing pitches. Different symbols are used for representing notes based on the duration of the note. All note symbols contain an ellipsoid that is either placed between two of the lines or touching one of the lines. Standard music notation may contain additional information such as \textit{accidentals}, different \textit{clefs}, agoic and several more that are not relevant in the context of this project. The following example shows the same melody notated in standard music notation.

\includegraphics[width= 0.4 \textwidth]{score_example}

Standard music notation will be referred to as \textit{Score} throughout the text.

\paragraph{Tablature}
Tablature is a simpler form of musical notation for string instruments. It has been developed prior to modern western musical notaion but has only gained a high degree of popularity throughout the last decades.

Tablature is based on a variable number of lines representing the string of the instrument. Numbers on the lines represent fret numbers.
Note that:
\begin{itemize}
	\item{Tablature typically does not feature any rhythmic information.}
	\item{On string instruments most notes can be played in more than one position of the fretboard. Tablature provides information on \textit{where} to play a certain note rathern than what note to play.}
\end{itemize}

The following example shows the same melody as above notated in tabulature for a six-string guitar.

\includegraphics[width= 0.4 \textwidth]{tab_example}

\section{Playing Techniques}
This section treats common playing techniques that are displayed, recognized and evaluated by the application. Emphasis is given to techniques relevant to string instruments.

\subsection{Transitions between Notes}
\paragraph{Bending}
\index{Bending}
Pulling a string upwards or downwards while playing a note causes the pitch to increase. In most situations strings are bent with the intention to raise the pitch of the note by a tone or a semitone.

\paragraph{Slide}
\index{Slide}
Moving the fretting finger up or down the neck of a string instrument while a note rings results in a smooth transition between two notes. This is referred to as \textit{sliding}.

\paragraph{Hammer On / Pull Off}
\index{Hammer On}
\index{Pull Off}
Similarly to slides, \textit{Hammer On}s and \textit{Pull Off}s are smooth transitions between notes. A Hammer On is executed by hammering on the fret, causing the note to sound without plucking the string with the right hand. Similarly a Pull Off is executed by pulling the string away from the fretboard while another finger is used for fretting the note.